\begin{definicion}[Tableros regulares y primera salida]\label{def:regular}
Fijamos $0<\delta<1$, $M=\lceil\delta n\rceil$ y un entero $K$ con $K_0\le K\le M$. El umbral $K_0$,
las constantes y las tolerancias que siguen se eligen en la Sección~\ref{meso:constantes}. Sea $B$ un
tablero de una historia, de tamaño $k$.
\begin{enumerate}[(a)]
\item \emph{Tramo macroscópico}, $k\ge M$. El tablero es \emph{regular} si la discrepancia $\Delta$
de la historia al llegar a él es menor que el radio $\varrho_{\rm mac}$ de la
Proposición~\ref{macro:parada}. La Sección~\ref{macro:observaciones} define esta discrepancia
con la partición $\mathcal I$, de paso $a_{\rm mac}$, de esa proposición.
\item \emph{Corte}, $k=M$. Se exige además que los observables de la
Sección~\ref{meso:regularidad} estén dentro de sus tolerancias de entrada y que
$D_{\max}(M)<\alpha_{\rm geo}M$, con la notación de la Sección~\ref{meso}.
\item \emph{Tramo mesoscópico}, $K\le k<M$. El tablero es regular si su punto de Bethe existe y
pertenece a la \emph{clase interior}: cumple las cotas de difusión, vacancia y brecha fijadas en la
Sección~\ref{meso:constantes} y la masa de cada diagonal es a lo sumo $\varepsilon$.
\end{enumerate}
La \emph{primera salida} de una historia es su primer tablero no regular, que llamamos \emph{tablero
de primera salida}. Su padre es regular. La historia se detiene en la primera salida o, si no la
hay, al llegar al tamaño $K$. El \emph{suceso de permanencia} $F$ consiste en que no haya salida:
todos los tableros hasta el de tamaño $K$, incluido éste, son regulares.
\end{definicion}
Para $n$ grande el tablero inicial es regular, pues su discrepancia es $O(1/n)$ (prueba de la
Proposición~\ref{macro:parada}). El punto de Bethe de un tablero regular cumple cotas de difusión,
vacancia y brecha con constantes fijas: en el tramo macroscópico por la Proposición~\ref{macro:maximo}
y el Corolario~\ref{macro:cambio}, y en el mesoscópico por definición. En el corte, las condiciones
adicionales lo sitúan ya en la clase interior, por la Proposición~\ref{meso:reconstruccion}. Sus
hijos, entre ellos el tablero de primera salida, heredan cotas del mismo tipo con constantes más
débiles, que llamamos \emph{cotas exteriores}; en el tramo mesoscópico, junto con una masa diagonal a
lo sumo $2\varepsilon$, definen la \emph{clase exterior}. En $\rho(B,i)$ intervienen todos los hijos
de $B$ en la fila $i$.

\begin{table}[htbp]\centering\small
\caption{Cómo se combinan las estimaciones en la comparación de los conteos.}
\begin{tabular}{@{}p{.31\linewidth}p{.59\linewidth}@{}}
\toprule Resultado & Estimaciones que utiliza \\ \midrule
Identidades \ref{hist:igualdad} y \ref{hist:inferior} & Partición por la siguiente fila y
normalización de las dos leyes.\\
Permanencia \ref{macro:cambio} & Balances \ref{macro:balances}, brecha \ref{macro:brecha},
reconstrucción \ref{macro:maximo} y parada \ref{macro:parada}.\\
Bandas \ref{prop:salida-banda-meso} & Prescripciones \eqref{prescr}, presupuesto
\ref{lema:presupuesto-meso}, concentración \ref{prop:concentracion-meso} y reconstrucción
\ref{meso:reconstruccion}.\\
Costes \ref{coste-macro} y \ref{coste-debil} & Comparación \ref{coste-comparacion} y, en el tramo
mesoscópico, normalizadores \ref{coste-normalizadores} y estimaciones del permanente
\ref{terminal-permanente} y \ref{terminal-normalizador}.\\
Último tablero \ref{terminal-error} & Geometría \ref{terminal-geometria}, pivotes
\ref{terminal-pivotes}, escalado \ref{terminal-escalado} y Proposición~3.1 de~\cite{dw}.\\
Teorema~\ref{princ} & Consistencia \ref{hist:consistencia}, cotas de salida y estimación
terminal bajo las dos leyes.\\ \bottomrule
\end{tabular}
\end{table}\begin{table}[!htbp]
\centering\small
\caption{Notación utilizada en la demostración.}
\label{tab:proof-notation}
\begin{tabular}{@{}p{.29\linewidth}>{\raggedright\arraybackslash}p{.65\linewidth}@{}}
\toprule
Símbolo & Significado \\
\midrule
$n,k$ & Orden del tablero inicial y número de filas vacías en el tablero residual. \\
$M,K$ & Corte $M=\lceil\delta n\rceil$ entre los dos tramos y tamaño terminal $K$. \\
$B_t,B^{(k)},B_\tau$ & Tablero tras $t$ inserciones, tablero con $k$ filas vacías
($B^{(k)}=B_{n-k}$) y tablero en la parada. \\
$\mathcal F_t,\mathcal F^{(k)}$ & Historia revelada, indexada por inserciones o por filas restantes,
con $\mathcal F^{(k)}=\mathcal F_{n-k}$. \\
\addlinespace[.22em]
$Q(B),L(B),g(B)$ & Conteo de completaciones, logaritmo del predictor y diferencia
$g(B)=\log Q(B)-L(B)$ cuando $Q(B)>0$. \\
$p_e,\ell_f,h_f$ & Peso de Bethe de la casilla $e$, masa de la diagonal $f$ y su vacancia
$h_f=1-\ell_f$. \\
\addlinespace[.22em]
$\mu,\nu,P^{\rm B}$ & Ley uniforme de revelado, ley adaptativa basada en el predictor
y ley auxiliar de Bethe. \\
$\tau,\rho(B,i),A_\tau$ & Tiempo de la primera salida o de la llegada al tamaño $K$, lo que ocurra
primero, cociente de normalización de la fila y logaritmo acumulado del cociente
$A_\tau=\sum_{t<\tau}\log\rho(B_t,i_t)$. \\
\addlinespace[.22em]
$F,G$ & Llegada al tamaño $K$ con regularidad hasta ese tablero inclusive, y subsuceso
$G=F\cap\{Q(B^{(K)})>0\}$. \\
$D_K$ & Suceso en el que dos cruces vacíos distintos en el tamaño $K$ comparten una diagonal. \\
\bottomrule
\end{tabular}
\end{table}

\FloatBarrier

\paragraph{Resultados que se usan de~\cite{pred}.}
Reunimos aquí, con una etiqueta, lo que las secciones siguientes toman de~\cite{pred}. Sea $B$ un
padre regular de tamaño $k$, con punto de Bethe $p$, sea $i$ una fila vacía y sea
$e\in R_i\cap E(B)$. En (I1)--(I3) reproducimos entre comillas la redacción de~\cite{pred}, donde la
casilla insertada se escribe $c$, su fila $r$, la corrección relativa $\delta_c$ y el resto $r_c$;
aquí escribimos $e$, $i$, $d_e$ y $r_e$.
\begin{enumerate}[({I}1)]
\item \emph{Herencia.} Para $k\ge k_0$, «cada hijo legal tiene un punto estacionario interior único
en la rama pequeña», es decir, con todos sus pesos menores que $1/4$, y «no se excluyen hijos
desfavorables ni se exige una cota inferior por casilla» \cite[Teorema~7.1]{pred}. En la clase
débil, «todos sus hijos legales tienen el punto estacionario pequeño construido en la
Sección~8.3» \cite[Teorema~8.1]{pred}. Sean $w$ la restricción de $p$ al soporte del hijo,
$\bar p=w+z$ su reequilibrado y $p^*$ el punto estacionario del hijo. Si $p$ cumple las cotas de
difusión, vacancia y brecha con constantes $C,\eta,\gamma$, para $k>12C/\gamma$ «todo hijo
restringido es conexo y su brecha es al menos $\gamma/2$»; además $|z_e|/w_e=O(k^{-1})$, las
vacancias de $\bar p$ son al menos $\eta/2$ y «la comparación multiplicativa de pesos conserva una
brecha uniforme» \cite[Sección~6.5]{pred}. Por último, $\|(p^*-\bar p)/\bar p\|_\infty=O(k^{-1})$ y
$p^*$ «permanece factible y positivo, con vacancias uniformes y $p^*<1/4$»
\cite[Sección~7.2]{pred}. En la clase débil valen las mismas cotas relativas y «la herencia es de
una clase interior a una exterior con margen» \cite[Secciones~8.2 y~8.3]{pred}. Los hijos legales son
todos los $B+e$ con $e\in E(B)$, tengan o no completaciones.
\item \emph{Desarrollo de un paso.} Se tiene
$|L(B+c)-L(B)-\log[p_c(1+\delta_c)]|\le C_*/k^2$ «para todo hijo $c$», con
$\delta_c=O(k^{-1})$ \cite[Teorema~7.1]{pred}. En la clase débil, $\delta_c=O(k^{-1})$ y
$r_c=L(B+c)-L(B)-\log[p_c(1+\delta_c)]=O(k^{-2})$ \cite[Secciones~7.8, 8.2 y~8.5]{pred}. Con
nuestra notación, $L(B+e)-L(B)=\log[p_e(1+d_e)]+r_e$, con $d_e=O(1/k)$ y $r_e=O(1/k^2)$
uniformemente en las casillas disponibles.
\item \emph{Cancelaciones.} En~\cite[Sección~6.1]{pred}, donde $B$ designa también la matriz de
las restricciones de filas y columnas, «$B(p\odot\delta)=BRg=0$» y, «en toda fila original,
$\sum_{c\in{\rm fila}}p_c\delta_c=0$»; además, «la restricción eliminada en cada componente se
reconstruye por la igualdad de las sumas de filas y columnas». Con nuestra notación, si
$B_{\rm RC}$ es esa matriz, $B_{\rm RC}(p\odot d)=0$: $\sum_{e'\in R_i}p_{e'}d_{e'}=0$ en cada fila
vacía y $\sum_{e'\in C_j}p_{e'}d_{e'}=0$ en cada columna vacía.
\item \emph{Consistencia.} $|\log\rho(B,i)|\le C_*/k^2$ \cite[Corolario~7.2 y Teorema~8.1]{pred}.
\item \emph{Constante del predictor.} $\log Q_{\rm pred}(n)-\log n!-n\Phi\to C_{\rm pred}$
\cite[Teoremas~1.1 y~5.7]{pred}.
\end{enumerate}
Las constantes implícitas, $k_0$ y $C_*$ dependen sólo de las cotas del padre y, cuando $k>M$,
también de $\delta$; en la clase débil no dependen de $n/k$. Los puntos (I4) e (I5), y las partes
de (I1) e (I2) citadas por un teorema, son enunciados numerados de~\cite{pred}; el resto de (I1)
e (I2) y el punto (I3) están en el texto de las secciones indicadas. Cada uso posterior repite el
lugar.
